% GENERATED FILE. Edit canonical lessons, then run npm run generate:books.
% canonical-source-hash: fnv1a64:d1e1823d26f18dde
% canonical-lessons: SA-C498-agrajah, SA-C498-anujah, SA-C498-jyesthah, SA-C498-kanisthah, SA-C498-yamalau

\chapter{Brothers and Sisters}
\label{ch:sa-brothers-and-sisters}
\hlchaptermodality{498}

\begin{chapteropening}
\textbf{By the end of this chapter:} \emph{I can say who is older and who is younger.}

Complete the last lesson of chapter 498: I can say who is older and who is younger.
\end{chapteropening}

\section[\texorpdfstring{\sk{अग्रजः} (agrajaḥ)}{अग्रजः (agrajaḥ)}]{\sk{अग्रजः} (agrajaḥ) — an elder brother}

\label{lesson:SA-C498-agrajah}



\begin{quote}
Before the new one: say the Sanskrit for offers into the fire, then the Sanskrit for good fortune.
\end{quote}

\subsection*{You'll want to know: \sk{अग्रजः}}
\textbf{\sk{अग्रजः}} — \emph{agrajaḥ} — \textquotedblleft{}an elder brother\textquotedblright{}.

\begin{quote}
\sk{मम} \sk{अग्रजः} \sk{नगरे} \sk{कार्यं} \sk{करोति।}
\end{quote}

\begin{quote}
\emph{mama agrajaḥ nagare kāryaṃ karoti.}
\end{quote}

\begin{quote}
My elder brother works in the city.
\end{quote}

A masculine noun.

\begin{cousinweb}
\textbf{\sk{अग्रजः}} (\emph{agrajaḥ}) is built from \emph{agra}, \textquotedblleft{}first, front\textquotedblright{}, and \emph{ja}, \textquotedblleft{}born\textquotedblright{}, from the root √jan; MW gives \textquotedblleft{}the first-born, an elder brother\textquotedblright{} (MW).
\end{cousinweb}

\begin{grammarlens}[title={What you've built}]
One more word for describing what you lived and felt.
\end{grammarlens}

\subsection*{Your turn}
\begin{itemize}
\raggedright
  \item \emph{Say it:} \emph{agrajaḥ}
  \item \emph{Say it:} \emph{agrajaḥ}, once more
  \item \emph{Say it:} \emph{mama agrajaḥ nagare kāryaṃ karoti.} — My elder brother works in the city.
  \item \emph{Recall:} say the Sanskrit for offers into the fire, then the Sanskrit for good fortune, then say \emph{agrajaḥ} again
\end{itemize}

\begin{tcolorbox}[breakable,colback=teal!4,colframe=teal!35!black,title={Before you move on}]
What does \sk{अग्रजः} mean? (\textquotedblleft{}An elder brother\textquotedblright{}.) Say it once more.
\end{tcolorbox}
\section[\texorpdfstring{\sk{अनुजः} (anujaḥ)}{अनुजः (anujaḥ)}]{\sk{अनुजः} (anujaḥ) — a younger brother}

\label{lesson:SA-C498-anujah}



\begin{quote}
Before the new one: say the Sanskrit for good fortune, then the Sanskrit for an elder brother.
\end{quote}

\subsection*{You'll want to know: \sk{अनुजः}}
\textbf{\sk{अनुजः}} — \emph{anujaḥ} — \textquotedblleft{}a younger brother\textquotedblright{}.

\begin{quote}
\sk{अनुजः} \sk{अग्रजेन} \sk{सह} \sk{क्रीडति।}
\end{quote}

\begin{quote}
\emph{anujaḥ agrajena saha krīḍati.}
\end{quote}

\begin{quote}
The younger brother plays with the elder.
\end{quote}

A masculine noun.

\begin{cousinweb}
\textbf{\sk{अनुजः}} (\emph{anujaḥ}) is built from \emph{anu}, \textquotedblleft{}after\textquotedblright{}, and \emph{ja}, \textquotedblleft{}born\textquotedblright{}, from the root √jan; MW gives \textquotedblleft{}born after, younger\textquotedblright{} and \textquotedblleft{}a younger brother\textquotedblright{} (MW).
\end{cousinweb}

\begin{grammarlens}[title={What you've built}]
One more word for describing what you lived and felt.
\end{grammarlens}

\subsection*{Your turn}
\begin{itemize}
\raggedright
  \item \emph{Say it:} \emph{anujaḥ}
  \item \emph{Say it:} \emph{anujaḥ}, once more
  \item \emph{Say it:} \emph{anujaḥ agrajena saha krīḍati.} — The younger brother plays with the elder.
  \item \emph{Recall:} say the Sanskrit for good fortune, then the Sanskrit for an elder brother, then say \emph{anujaḥ} again
\end{itemize}

\begin{tcolorbox}[breakable,colback=teal!4,colframe=teal!35!black,title={Before you move on}]
What does \sk{अनुजः} mean? (\textquotedblleft{}A younger brother\textquotedblright{}.) Say it once more.
\end{tcolorbox}
\section[\texorpdfstring{\sk{ज्येष्ठः} (jyeṣṭhaḥ)}{ज्येष्ठः (jyeṣṭhaḥ)}]{\sk{ज्येष्ठः} (jyeṣṭhaḥ) — eldest}

\label{lesson:SA-C498-jyesthah}



\begin{quote}
Before the new one: say the Sanskrit for an elder brother, then the Sanskrit for a younger brother.
\end{quote}

\subsection*{You'll want to know: \sk{ज्येष्ठः}}
\textbf{\sk{ज्येष्ठः}} — \emph{jyeṣṭhaḥ} — \textquotedblleft{}eldest\textquotedblright{}.

\begin{quote}
\sk{ज्येष्ठः} \sk{पुत्रः} \sk{पितुः} \sk{साहाय्यं} \sk{करोति।}
\end{quote}

\begin{quote}
\emph{jyeṣṭhaḥ putraḥ pituḥ sāhāyyaṃ karoti.}
\end{quote}

\begin{quote}
The eldest son helps his father.
\end{quote}

\begin{cousinweb}
\textbf{\sk{ज्येष्ठः}} (\emph{jyeṣṭhaḥ}) MW gives it as a superlative (Pāṇini 5.3.61), \textquotedblleft{}pre-eminent, greatest, eldest\textquotedblright{}, without naming a root, so its origin is uncertain (MW).
\end{cousinweb}

\begin{grammarlens}[title={What you've built}]
One more word for describing what you lived and felt.
\end{grammarlens}

\subsection*{Your turn}
\begin{itemize}
\raggedright
  \item \emph{Say it:} \emph{jyeṣṭhaḥ}
  \item \emph{Say it:} \emph{jyeṣṭhaḥ}, once more
  \item \emph{Say it:} \emph{jyeṣṭhaḥ putraḥ pituḥ sāhāyyaṃ karoti.} — The eldest son helps his father.
  \item \emph{Recall:} say the Sanskrit for an elder brother, then the Sanskrit for a younger brother, then say \emph{jyeṣṭhaḥ} again
\end{itemize}

\begin{tcolorbox}[breakable,colback=teal!4,colframe=teal!35!black,title={Before you move on}]
What does \sk{ज्येष्ठः} mean? (\textquotedblleft{}Eldest\textquotedblright{}.) Say it once more.
\end{tcolorbox}
\section[\texorpdfstring{\sk{कनिष्ठः} (kaniṣṭhaḥ)}{कनिष्ठः (kaniṣṭhaḥ)}]{\sk{कनिष्ठः} (kaniṣṭhaḥ) — youngest}

\label{lesson:SA-C498-kanisthah}



\begin{quote}
Before the new one: say the Sanskrit for a younger brother, then the Sanskrit for eldest.
\end{quote}

\subsection*{You'll want to know: \sk{कनिष्ठः}}
\textbf{\sk{कनिष्ठः}} — \emph{kaniṣṭhaḥ} — \textquotedblleft{}youngest\textquotedblright{}.

\begin{quote}
\sk{कनिष्ठा} \sk{भगिनी} \sk{मम} \sk{प्रिया} \sk{अस्ति।}
\end{quote}

\begin{quote}
\emph{kaniṣṭhā bhaginī mama priyā asti.}
\end{quote}

\begin{quote}
My youngest sister is my favourite.
\end{quote}

\begin{cousinweb}
\textbf{\sk{कनिष्ठः}} (\emph{kaniṣṭhaḥ}) MW gives \textquotedblleft{}the youngest, younger born\textquotedblright{}, opposed to \emph{jyeṣṭha}, without naming a root, so its origin is uncertain (MW).
\end{cousinweb}

\begin{grammarlens}[title={What you've built}]
One more word for describing what you lived and felt.
\end{grammarlens}

\subsection*{Your turn}
\begin{itemize}
\raggedright
  \item \emph{Say it:} \emph{kaniṣṭhaḥ}
  \item \emph{Say it:} \emph{kaniṣṭhaḥ}, once more
  \item \emph{Say it:} \emph{kaniṣṭhā bhaginī mama priyā asti.} — My youngest sister is my favourite.
  \item \emph{Recall:} say the Sanskrit for a younger brother, then the Sanskrit for eldest, then say \emph{kaniṣṭhaḥ} again
\end{itemize}

\begin{tcolorbox}[breakable,colback=teal!4,colframe=teal!35!black,title={Before you move on}]
What does \sk{कनिष्ठः} mean? (\textquotedblleft{}Youngest\textquotedblright{}.) Say it once more.
\end{tcolorbox}
\section[\texorpdfstring{\sk{यमलौ} (yamalau)}{यमलौ (yamalau)}]{\sk{यमलौ} (yamalau) — twins}

\label{lesson:SA-C498-yamalau}



\begin{quote}
Before the new one: say the Sanskrit for eldest, then the Sanskrit for youngest.
\end{quote}

\subsection*{You'll want to know: \sk{यमलौ}}
\textbf{\sk{यमलौ}} — \emph{yamalau} — \textquotedblleft{}twins\textquotedblright{}.

\begin{quote}
\sk{यमलौ} \sk{बालौ} \sk{सह} \sk{क्रीडतः।}
\end{quote}

\begin{quote}
\emph{yamalau bālau saha krīḍataḥ.}
\end{quote}

\begin{quote}
The twin boys play together.
\end{quote}

A masculine noun, used in the dual.

\begin{cousinweb}
\textbf{\sk{यमलौ}} (\emph{yamalau}) is built on \emph{yama}, \textquotedblleft{}a twin\textquotedblright{}; MW gives \textquotedblleft{}twin, paired\textquotedblright{} and, in the dual, \textquotedblleft{}twins, a pair\textquotedblright{} (MW).
\end{cousinweb}

\begin{grammarlens}[title={What you've built}]
One more word for describing what you lived and felt.
\end{grammarlens}

\subsection*{Your turn}
\begin{itemize}
\raggedright
  \item \emph{Say it:} \emph{yamalau}
  \item \emph{Say it:} \emph{yamalau}, once more
  \item \emph{Say it:} \emph{yamalau bālau saha krīḍataḥ.} — The twin boys play together.
  \item \emph{Recall:} say the Sanskrit for eldest, then the Sanskrit for youngest, then say \emph{yamalau} again
\end{itemize}

\begin{tcolorbox}[breakable,colback=teal!4,colframe=teal!35!black,title={Before you move on}]
What does \sk{यमलौ} mean? (\textquotedblleft{}Twins\textquotedblright{}.) Say it once more.
\end{tcolorbox}
